\documentclass[10pt,a4paper]{amsart}
\usepackage{amsmath,amssymb,amsthm,graphicx}
\usepackage{verbatim}
\newenvironment{reference}{\comment}{\endcomment}
%\newenvironment{reference}{}{}
\newenvironment{slogan}{\comment}{\endcomment}
\newenvironment{history}{\comment}{\endcomment}

% For commutative diagrams we use Xy-pic
\usepackage[all]{xy}

% We use 2cell for 2-commutative diagrams.
\xyoption{2cell}
% Theorem environments.
%
\theoremstyle{plain}
\newtheorem{theorem}[subsection]{Theorem}
\newtheorem{proposition}[subsection]{Proposition}
\newtheorem{lemma}[subsection]{Lemma}

\theoremstyle{definition}
\newtheorem{definition}[subsection]{Definition}
\newtheorem{example}[subsection]{Example}
\newtheorem{exercise}[subsection]{Exercise}
\newtheorem{situation}[subsection]{Situation}

\theoremstyle{remark}
\newtheorem{remark}[subsection]{Remark}
\newtheorem{remarks}[subsection]{Remarks}

\numberwithin{equation}{subsection}

% Macros
%
\def\lim{\mathop{\mathrm{lim}}\nolimits}
\def\colim{\mathop{\mathrm{colim}}\nolimits}
\def\Spec{\mathop{\mathrm{Spec}}}
\def\Hom{\mathop{\mathrm{Hom}}\nolimits}
\def\Ext{\mathop{\mathrm{Ext}}\nolimits}
\def\SheafHom{\mathop{\mathcal{H}\!\mathit{om}}\nolimits}
\def\SheafExt{\mathop{\mathcal{E}\!\mathit{xt}}\nolimits}
\def\Sch{\mathit{Sch}}
\def\Mor{\mathop{\mathrm{Mor}}\nolimits}
\def\Ob{\mathop{\mathrm{Ob}}\nolimits}
\def\Sh{\mathop{\mathit{Sh}}\nolimits}
\def\NL{\mathop{N\!L}\nolimits}
\def\CH{\mathop{\mathrm{CH}}\nolimits}
\def\proetale{{pro\text{-}\acute{e}tale}}
\def\etale{{\acute{e}tale}}
\def\QCoh{\mathit{QCoh}}
\def\Ker{\mathop{\mathrm{Ker}}}
\def\Im{\mathop{\mathrm{Im}}}
\def\Coker{\mathop{\mathrm{Coker}}}
\def\Coim{\mathop{\mathrm{Coim}}}

% Boxtimes
%
\DeclareMathSymbol{\boxtimes}{\mathbin}{AMSa}{"02}

%
% Macros for moduli stacks/spaces
%
\def\QCohstack{\mathcal{QC}\!\mathit{oh}}
\def\Cohstack{\mathcal{C}\!\mathit{oh}}
\def\Spacesstack{\mathcal{S}\!\mathit{paces}}
\def\Quotfunctor{\mathrm{Quot}}
\def\Hilbfunctor{\mathrm{Hilb}}
\def\Curvesstack{\mathcal{C}\!\mathit{urves}}
\def\Polarizedstack{\mathcal{P}\!\mathit{olarized}}
\def\Complexesstack{\mathcal{C}\!\mathit{omplexes}}
% \Pic is the operator that assigns to X its picard group, usage \Pic(X)
% \Picardstack_{X/B} denotes the Picard stack of X over B
% \Picardfunctor_{X/B} denotes the Picard functor of X over B
\def\Pic{\mathop{\mathrm{Pic}}\nolimits}
\def\Picardstack{\mathcal{P}\!\mathit{ic}}
\def\Picardfunctor{\mathrm{Pic}}
\def\Deformationcategory{\mathcal{D}\!\mathit{ef}}
\usepackage{cleveref}
\crefname{originallemma}{original lemma}{original lemmas}
\Crefname{originallemma}{Original lemma}{Original lemmas}
\makeatletter
\newlabel{lemma-localize-given-products}{{03CE}{source}{}{sourcepointer-1}{}}
\newlabel{lemma-localize-given-products@cref}{{[originallemma][0][]03CE}{[1][0][]source}{}{}{}}
\makeatother
\crefname{originallemma}{original lemma}{original lemmas}
\Crefname{originallemma}{Original lemma}{Original lemmas}
\makeatletter
\newlabel{lemma-localize-morphism}{{03CF}{source}{}{sourcepointer-2}{}}
\newlabel{lemma-localize-morphism@cref}{{[originallemma][0][]03CF}{[1][0][]source}{}{}{}}
\makeatother
\crefname{originalproposition}{original proposition}{original propositions}
\Crefname{originalproposition}{Original proposition}{Original propositions}
\makeatletter
\newlabel{proposition-get-morphism}{{00X6}{source}{}{sourcepointer-3}{}}
\newlabel{proposition-get-morphism@cref}{{[originalproposition][0][]00X6}{[1][0][]source}{}{}{}}
\makeatother
\crefname{originaldefinition}{original definition}{original definitions}
\Crefname{originaldefinition}{Original definition}{Original definitions}
\makeatletter
\newlabel{definition-composition-morphisms-sites}{{03CC}{source}{}{sourcepointer-4}{}}
\newlabel{definition-composition-morphisms-sites@cref}{{[originaldefinition][0][]03CC}{[1][0][]source}{}{}{}}
\makeatother
\title{Localizing a morphism of sites when the functor commutes with finite limits}
\author{The Stacks Project Authors}
\date{stacks-project:sites.tex:local-snapshot-20261012; GNU FDL 1.2 or later}
\begin{document}
\maketitle
\noindent\textit{Editorial note.} All delivered mathematical text is the author's own, copied from the pinned Stacks Project chapter file sites-cohomology.tex; no statement or proof step was rewritten, reordered or removed.\par\medskip
\noindent\textit{Editorial note.} The proof is the author's complete argument; every step the author writes out is inside the counted span and no closing sentence is a placeholder.\par\medskip
\noindent\textit{Editorial note.} Every labelled reference occurring inside the counted spans is carried in assembly.external\_references as a non-load-bearing navigation pointer (same-chapter ones at their definition site, cross-chapter ones at the invocation inside the counted span), so no load-bearing mathematics is replaced by a stub. No bibliography citation occurs inside the counted spans.\par\medskip
\noindent\textit{Editorial note.} The checked chapter list for same-name collisions is the set of stacks chapters staged in this repository: chow, cohomology, divisors, etale-cohomology, flat, more-algebra, more-morphisms, morphisms, spaces-more-morphisms, varieties together with this card's chapter sites-cohomology.\par\medskip
\noindent\textit{Editorial note.} The preamble is the chapter's own retained minimal preamble (sources/stacks-sites-cohomology/retained-preamble.tex), a verbatim excerpt of the native preamble: the Xy-pic setup and the native theorem declarations with the mathematical macros.\par\medskip
\noindent\textit{Allowed standard prerequisites (editorial).}\begin{itemize}
\item The proposition producing a morphism of sites from a functor. Same chapter, Proposition proposition-get-morphism. Produces the two morphisms of sites appearing in the diagram.
\item The localization of a site when the site has products. Same chapter, Lemma lemma-localize-given-products. Supplies the description of the horizontal arrows of the diagram.
\item Composition of morphisms of sites. Same chapter, Definition definition-composition-morphisms-sites. Supplies the commutativity of the diagram of categories and functors.
\end{itemize}\medskip
\ifdefined\hypertarget\hypertarget{sourcepointer-1}{}\fi
\noindent\textit{Source reference.} Original lemma 03CE: sources/stacks-sites/pinned-original-source.tex, source lines 6629--6629. Navigation pointer invoked at proof line 6629; the counted span contains the pointer only and the cited result is not reproduced inside it, so the pointer is registered as a non-load-bearing navigation stub in the established form for Stacks material. No bibliography citation is involved.\par\medskip
\ifdefined\hypertarget\hypertarget{sourcepointer-2}{}\fi
\noindent\textit{Source reference.} Original lemma 03CF: sources/stacks-sites/pinned-original-source.tex, source lines 6601--6601. Navigation pointer invoked at proof line 6601; the counted span contains the pointer only and the cited result is not reproduced inside it, so the pointer is registered as a non-load-bearing navigation stub in the established form for Stacks material. No bibliography citation is involved.\par\medskip
\ifdefined\hypertarget\hypertarget{sourcepointer-3}{}\fi
\noindent\textit{Source reference.} Original proposition 00X6: sources/stacks-sites/pinned-original-source.tex, source lines 6638--6638. Navigation pointer invoked at proof line 6638; the counted span contains the pointer only and the cited result is not reproduced inside it, so the pointer is registered as a non-load-bearing navigation stub in the established form for Stacks material. No bibliography citation is involved.\par\medskip
\ifdefined\hypertarget\hypertarget{sourcepointer-4}{}\fi
\noindent\textit{Source reference.} Original definition 03CC: sources/stacks-sites/pinned-original-source.tex, source lines 6645--6645. Navigation pointer invoked at proof line 6645; the counted span contains the pointer only and the cited result is not reproduced inside it, so the pointer is registered as a non-load-bearing navigation stub in the established form for Stacks material. No bibliography citation is involved.\par\medskip
\begin{lemma}
\label{lemma-localize-morphism-strong}
Let $\mathcal{C}$, $\mathcal{D}$ be sites.
Let $u : \mathcal{D} \to \mathcal{C}$ be a functor.
Let $V \in \Ob(\mathcal{D})$. Set $U = u(V)$.
Assume that
\begin{enumerate}
\item $\mathcal{C}$ and $\mathcal{D}$ have
all finite limits,
\item $u$ is continuous, and
\item $u$ commutes with finite limits.
\end{enumerate}
There exists a commutative diagram of morphisms of sites
$$
\xymatrix{
\mathcal{C}/U \ar[r]_{j_U} \ar[d]_{f'} & \mathcal{C} \ar[d]^f \\
\mathcal{D}/V \ar[r]^{j_V} & \mathcal{D}
}
$$
where the right vertical arrow corresponds to $u$,
the left vertical arrow corresponds to the
functor $u' : \mathcal{D}/V \to \mathcal{C}/U$, $V'/V \mapsto u(V')/u(V)$
and the horizontal arrows correspond to the functors
$\mathcal{C} \to \mathcal{C}/U$, $X \mapsto X \times U$
and $\mathcal{D} \to \mathcal{D}/V$, $Y \mapsto Y \times V$
as in Lemma \ref{lemma-localize-given-products}.
Moreover, the associated diagram of morphisms of topoi is
equal to the diagram of
Lemma \ref{lemma-localize-morphism}.
In particular we have $f'_*j_U^{-1} = j_V^{-1}f_*$.
\end{lemma}

\begin{proof}
Note that $u$ satisfies the assumptions of
Proposition \ref{proposition-get-morphism} and hence induces
a morphism of sites $f : \mathcal{C} \to \mathcal{D}$ by that proposition.
It is clear that $u$ induces a functor $u'$ as indicated.
It is clear that this functor also satisfies the assumptions of
Proposition \ref{proposition-get-morphism}.
Hence we get a morphism of sites $f' : \mathcal{C}/U \to \mathcal{D}/V$.
The diagram commutes by our definition of composition of morphisms of
sites (see Definition \ref{definition-composition-morphisms-sites})
and because
$$
u(Y \times V) = u(Y) \times u(V) = u(Y) \times U
$$
which shows that the diagram of categories and functors opposite to
the diagram of the lemma commutes.
\end{proof}

\end{document}
